\scanpage{102}
Ce qui signifie aussi que tt \inserted{(ou un)} scindage
de 31 --- p.ex. défini par un pt de $U_\eta$
\unsure{rat|rt}/$K$ --- induit un hom $T\longrightarrow\pi$)

\noindent c) L'action de $T$ sur $\pi_{\mathrm{ab}}=H_1(U_{\bar\eta})$ est
triviale.

\noindent d) Item pour $H_{\unsure{1|n}}(U_{\bar\eta},T_\ell)$, si $\ell$ est premier
à la car. résiduelle. (\uline{NB} \deleted{\unsure{Pour car. $p>0$|Pour $\ell\ne p$|Pour $\ell=p$}} \inserted{$\ell=p>0$}\unsure{,|)}
on doit pouvoir l'exprimer par un critère
sur le groupe $p$-divisible\ldots)

\noindent e) (En termes d'une section de (30),
définissant une extension de $\mathcal E_k$ par $\pi$,
et supposant maintenant $k$ de type
fini sur $\mathbf Q$) La dite extension est
admissible (quitte à passer à un
sous groupe ouvert \inserted{de $\mathcal E_k$,} au besoin), ie.
peut être défini par une courbe alg.
sur \inserted{$k$, ou sur} une extension finie de $k$.

\noindent f) (En termes d'une section de (30)
et d'une de (31)., permettant de
faire opérer $\mathcal E_k$ sur $\pi$) : Cette
action \inserted{n'}est \inserted{pas trop} loufoque (à définir convenablement !)% The end of convenablement continues on a short lower baseline; no hyphen is written.
\margin{101}
